%==============================================================
%  Paper II — fractal tumor-to-brain descriptor (Pillar I)
%  IEEEtran journal format. Compile on Overleaf with pdfLaTeX.
%==============================================================
\documentclass[journal]{IEEEtran}

\usepackage{cite}
\usepackage{amsmath,amssymb,amsfonts}
\usepackage{algorithmic}
\usepackage{algorithm}
\usepackage{graphicx}
\usepackage{textcomp}
\usepackage{booktabs}
\usepackage{multirow}
\usepackage{url}
\hyphenation{me-nin-gi-o-ma}

\begin{document}

\title{A Tumor-to-Brain Fractal Ratio for Interpretable\\
Brain Tumor Classification in MRI:\\
Analysis, Validation, and Limits}

\author{Zahraa Ibrahim Khaleel and Mohammed Sahib Mehdi%
\thanks{Manuscript received \today.}
\thanks{The authors are with the Department of Computer Science, College of Science, Al-Nahrain University, Baghdad, Iraq (e-mail: [YOUR EMAIL]).}}

\markboth{Journal manuscript,~Vol.~XX, No.~X, 2026}%
{Khaleel \MakeLowercase{\textit{and}} Mehdi: A Tumor-to-Brain Fractal Ratio for Interpretable Brain Tumor Classification}

\maketitle

\begin{abstract}
Deep networks classify brain tumor type from MRI with high reported accuracy,
but the features they use cannot be named, and the accuracies most often quoted
are not measured at the level at which a clinical decision is made. We take the
opposite position and ask how far a small set of named, measurable quantities can
go. The central quantity is a ratio, $C = F_t/F_b$, between the box-counting
fractal dimension of the tumor region and that of the same patient's brain with
the tumor removed. Because the two regions come from one acquisition, the ratio
cancels the scanner and sequence factors that make a raw fractal dimension
incomparable between patients, and it expresses how irregular a tumor is relative
to the brain that hosts it. Combined with fifteen further interpretable
descriptors of texture, entropy, lacunarity, shape and edge sharpness, and
classified with a random forest under patient-grouped cross-validation, the set
reaches a patient-level balanced accuracy of $0.92$ on the Figshare brain tumor
dataset, against $0.913$ for the benchmark published with that dataset. We also
report what the descriptor does not do. Taken alone it separates pituitary
adenomas well (AUC $0.907$) and meningiomas barely (AUC $0.613$), and $45\%$ of
its class separation survives the removal of tumor size, so it is neither a pure
morphology measure nor a proxy for size. Two published methods obtain higher
accuracy on the same data using learned or higher-dimensional representations.
The contribution is therefore a characterisation: what a single interpretable
ratio measures, how much of the task it can carry, and where it fails.
\end{abstract}

\begin{IEEEkeywords}
Brain tumor classification, fractal dimension, box counting, interpretable
features, magnetic resonance imaging, patient-level evaluation.
\end{IEEEkeywords}

%==============================================================
\section{Introduction}
\IEEEPARstart{T}{he} three most common primary intracranial tumors seen on
contrast-enhanced T1-weighted MRI, meningioma, glioma and pituitary adenoma,
differ in ways a radiologist describes verbally: one is well circumscribed and
attached to the dura, another infiltrates surrounding parenchyma with an
irregular margin, the third occupies a characteristic anatomical location. These
descriptions are qualitative, and the interest in automating them is long
standing \cite{cheng2015,gumaei2019,deepak2019}.

Most recent progress on this task has come from transferring convolutional
networks trained on natural images \cite{deepak2019,swati2019} or from
architectures designed to preserve spatial relationships \cite{afshar2019}.
Reported accuracies on the standard benchmark now exceed $95\%$. Two properties
of that literature motivate the present work.

The first is that the features are not nameable. A network that reaches high
accuracy provides a decision but not a quantity a clinician can inspect, compare
against a measurement made by hand, or reason about when it is wrong. Saliency
maps indicate where a network attended, which is not the same as stating what it
measured. In settings where a prediction must be justified rather than merely
supplied, this matters.

The second is that reported accuracies are frequently not comparable to each
other. Each patient in this dataset contributes many adjacent slices with almost
identical appearance. A random split over slices therefore places near-duplicates
in both training and test sets. Yagis \emph{et al.} measured the size of this
effect on brain MRI and found that slice-level splitting inflated accuracy by
$29$ to $55$ points relative to subject-level splitting, and that scans with
randomly assigned labels reached approximately $96\%$ accuracy under a
slice-level split against $50\%$ under a subject-level split \cite{yagis2021}.
Those measurements were made on Alzheimer's and Parkinson's cohorts rather than
on tumor data, but the mechanism is a property of how the images are organised
and not of the pathology: wherever one subject contributes many adjacent slices,
the same inflation is available. A number produced under the two protocols is
not the same number.

This paper asks what a small set of quantities that can each be named and
measured achieves under an evaluation protocol that does not inflate itself. The
descriptor at the centre of it is a ratio between two fractal dimensions computed
on the same slice.

Fractal analysis of brain tissue is not new. Box-counting dimensions were applied
to tumor detection in MRI by Iftekharuddin \emph{et al.} \cite{iftekhar2003}.
Dimensions computed over both tumorous and non-tumorous tissue have been used for
glioma grading \cite{frontiers2023}, and fractal measures combined with a
boundary regularity ratio have been used to separate tumor types
\cite{martin2022}. What we have not found in that literature is the use of the
single ratio $F_t/F_b$, tumor dimension over tumor-free brain dimension, as a
descriptor for three-class tumor typing, and the analysis of what such a ratio
does and does not capture.

The contributions are as follows.

\begin{itemize}
\item We define $C = F_t/F_b$ and give the reason it is comparable across
patients where a raw fractal dimension is not: both regions come from one
acquisition, so the multiplicative factors introduced by scanner, sequence and
noise divide out.

\item We measure what the descriptor captures. Its one-vs-rest AUC is $0.907$
for pituitary adenoma, $0.735$ for glioma and $0.613$ for meningioma, so its
strength is strongly class dependent, and we explain why the morphological
heterogeneity of meningioma is the likely reason.

\item We test the descriptor against the obvious alternative explanation. After
removing the component predictable from tumor area, $45\%$ of its class
separation remains, so it is partly but not wholly a proxy for size.

\item We report classification with sixteen interpretable features under
patient-grouped cross-validation, reaching $0.92$ balanced accuracy at the
patient level, and we position that figure honestly against published results,
two of which are higher.
\end{itemize}

%==============================================================
\section{Related Work}

\subsection{Fractal Analysis of Brain Tissue}
The fractal dimension quantifies how the detail in a structure changes with the
scale at which it is measured, and tissue that is structurally disorganised
yields a higher value than tissue that is regular. Iftekharuddin \emph{et al.}
introduced box-counting variants for tumor detection and localisation in brain
MRI \cite{iftekhar2003}. More recent work computes several fractal measures over
tumorous and non-tumorous regions separately for glioma grading, reporting that
measures taken over apparently normal tissue also carry grading information
\cite{frontiers2023}. Sánchez and Martín-Landrove combine the dimension of the
tumor interface with a regularity ratio to separate gliomas from meningiomas and
schwannomas \cite{martin2022}, and treat lacunarity as a complementary measure of
how gappy a region is at a given scale.

These works establish that fractal measures carry information about brain tumors.
They compute the measures as separate features. The ratio used here is a
normalisation, and its purpose is comparability across patients rather than the
addition of a new measurement.

\subsection{Classification on the Figshare Benchmark}
The dataset published by Cheng \emph{et al.} \cite{cheng2015} is the standard
reference for three-class tumor typing, and the accompanying method, which
augments the tumor region and partitions it into rings before a bag-of-words
representation, reports $91.28\%$.

Later work follows two directions. Engineered representations include the
discrete wavelet and Gabor statistics of Ismael and Abdel-Qader
\cite{ismael2018}, at $91.9\%$, and the hybrid PCA-normalised GIST features with
a regularised extreme learning machine of Gumaei \emph{et al.}
\cite{gumaei2019}, at $94.23\%$. Learned representations include the GoogLeNet
features with a support vector machine of Deepak and Ameer \cite{deepak2019}, at
$98\%$ under patient-level five-fold cross-validation, the block-wise fine-tuned
VGG19 of Swati \emph{et al.} \cite{swati2019} at $94.82\%$, and capsule networks
operating on segmented tumors and coarse boundaries \cite{afshar2019} at
$90.98\%$.

We take these numbers at face value where the evaluation protocol is stated, and
we note where it is not, since the difference between protocols is larger than
the differences between many of the methods.

\subsection{Interpretability as a Design Constraint}
The argument for named features in clinical imaging is not that they are more
accurate. It is that a quantity with a name can be audited. If a model assigns a
lesion to a class, a clinician can ask which measurement drove the assignment,
recompute that measurement from the image, and disagree with it on grounds that
can be stated. A saliency map indicates the region a network attended to but does
not supply a quantity, so the same audit is not available.

This distinction matters more when the training set is small. With $233$ patients
a learned representation has many more free parameters than the data constrain,
and the features it settles on need not correspond to anything stable. A fixed
set of measurements has no such freedom: the value of lacunarity for a given mask
is the same regardless of what the training set contained.

The counter-argument is equally clear and is not disputed here. Engineered
features encode a prior about what matters, and where that prior is wrong the
representation cannot recover. The results below show one instance of exactly
that, in the class for which the descriptor is weakest.

\subsection{Evaluation Protocol}
The distinction between slice-level and patient-level evaluation is treated here
as a methodological requirement rather than a detail. The quantitative basis is
the leakage study of Yagis \emph{et al.} \cite{yagis2021} described in the
introduction. We group by patient identifier throughout and report patient-level
figures, and we compare only against studies that state a patient-respecting
protocol.

%==============================================================
\section{Methods}

\subsection{Region Separation}
Three regions are derived from each slice, and the descriptor is defined over
two of them.

The \emph{brain region} is obtained by thresholding the slice at
$\max(0.25\bar{I}, 10)$ where $\bar{I}$ is the mean intensity, filling holes,
applying a binary opening with two iterations, and retaining the largest
connected component. This step exists to exclude the black background, which
would otherwise contribute empty boxes to every box-counting estimate and bias
the dimension downward by an amount that depends on how much of the field of view
the head occupies.

The \emph{tumor region} is the expert mask supplied with the dataset,
intersected with the brain region so that no mask pixel lying outside the
estimated brain contributes.

The \emph{tumor-free brain} is the brain region with the tumor region removed.

We use the expert masks rather than an automatic segmentation because the
descriptor is the object under study. Computing it on an automatically produced
mask would confound the properties of the descriptor with the errors of the
segmenter, and a change in either would be indistinguishable in the result. The
consequence, which we state rather than set aside, is that the method as
evaluated requires a mask at inference time; Section~\ref{sec:limits} returns to
this.

A slice is excluded if the tumor region contains fewer than $16$ pixels or the
tumor-free region fewer than $64$, since a box-counting estimate over a region
that small is not meaningful. On the Figshare dataset this leaves $2740$ of
$3064$ slices, from $218$ of $233$ patients.

\subsection{The Fractal Descriptor}
For a region $R$ within a slice, intensities inside $R$ are rescaled to
$[0,255]$ using the minimum and maximum within $R$, so that the estimate depends
on the structure of $R$ rather than on the global intensity range of the slice.
The plane is covered with boxes of side $s$, and for each box intersecting $R$
the number of intensity levels spanned within that box is counted, following
differential box counting. Writing $N(s)$ for the total count at scale $s$, the
dimension is the slope of the least-squares fit
\begin{equation}
\log N(s) \;=\; F \cdot \log\left(\tfrac{1}{s}\right) + b ,
\qquad s \in \{2, 4, \ldots, 128\}.
\label{eq:fd}
\end{equation}

The descriptor is the ratio of the dimension of the tumor region to that of the
tumor-free brain in the same slice,
\begin{equation}
C \;=\; \frac{F_t}{F_b}.
\label{eq:C}
\end{equation}

The reason for the ratio is comparability. A raw fractal dimension depends on
acquisition parameters, field strength, contrast administration and noise, all of
which vary between patients and institutions, so $F_t$ measured on one patient is
not directly comparable to $F_t$ measured on another. The tumor-free brain in the
same slice was acquired under exactly the same conditions. If those conditions
act as a common multiplicative factor on both regions, the factor divides out and
$C$ expresses how irregular the tumor is \emph{relative to the brain that hosts
it}. On the Figshare dataset $C$ takes values in $[0.428, 1.076]$.

\begin{algorithm}[t]
\caption{Computation of the tumor-to-brain fractal ratio}
\label{alg:c}
\begin{algorithmic}[1]
\STATE \textbf{Input:} slice $I$, expert tumor mask $M$
\STATE $B \leftarrow$ threshold $I$ at $\max(0.25\bar{I},10)$, fill holes, open, keep largest component
\STATE $T \leftarrow M \cap B$ \COMMENT{tumor region}
\STATE $B_0 \leftarrow B \setminus T$ \COMMENT{tumor-free brain}
\IF{$|T| < 16$ \OR $|B_0| < 64$}
  \STATE \textbf{reject slice}
\ENDIF
\FOR{$R \in \{T, B_0\}$}
  \STATE rescale intensities within $R$ to $[0,255]$
  \FOR{$s \in \{2,4,\ldots,128\}$}
    \STATE $N_R(s) \leftarrow \sum_{\text{boxes } b \cap R \neq \emptyset}
            \left( \lfloor \max_b / s \rfloor - \lfloor \min_b / s \rfloor + 1 \right)$
  \ENDFOR
  \STATE $F_R \leftarrow$ slope of $\log N_R(s)$ against $\log(1/s)$
\ENDFOR
\STATE \textbf{return} $C = F_T / F_{B_0}$
\end{algorithmic}
\end{algorithm}

Algorithm~\ref{alg:c} states the computation. Two details affect the value and
are fixed for all reported results. The box sizes are powers of two up to the
image side, so the fit in \eqref{eq:fd} is over eight points evenly spaced in
$\log s$; using fewer points makes the slope estimate unstable on small regions,
and using non-dyadic sizes introduces partial boxes at the region boundary whose
treatment would need a further convention. The intensity rescaling is performed
within each region separately rather than over the slice, so that $F_t$ describes
the structure of the tumor rather than its intensity relative to the rest of the
image; without this, a tumor that is uniformly bright would receive a different
dimension from an identical tumor in a darker slice.

\subsection{The Remaining Descriptors}
The fifteen features accompanying $C$ are computed as follows, and each is
reported here in enough detail to be recomputed.

\emph{Texture.} A grey-level co-occurrence matrix \cite{haralick1973} is formed
at unit offset and zero angle on intensities quantised to eight levels, computed
separately over the tumor and the tumor-free brain. Contrast is taken from both
regions and correlation from the tumor. Eight levels rather than the conventional
sixteen or thirty-two are used because the tumor regions are small, and a finer
quantisation leaves most matrix entries empty.

\emph{Entropy.} The Shannon entropy of tumor intensities over a $32$-bin
histogram, which measures heterogeneity of intensity independently of its spatial
arrangement, and so complements the co-occurrence features.

\emph{Lacunarity.} Following the gliding-box definition \cite{allain1991}, the
mask is covered with boxes of side eight, the mask pixel count in each box is
recorded, and lacunarity is $1 + \sigma^2/\mu^2$ over those counts. A region can
share a fractal dimension with another while differing in lacunarity, because the
first describes how detail scales and the second how the region is distributed at
one scale.

\emph{Shape.} Solidity, eccentricity, extent and circularity
$4\pi A / P^2$ are computed on the largest connected component of the tumor mask.
These describe the outline rather than the interior and are the features expected
to separate a well-circumscribed lesion from an infiltrating one.

\emph{Edge.} The intensity gradient magnitude is computed by finite differences.
Edge density is the fraction of tumor pixels whose gradient exceeds the slice
mean; mean gradient is taken over the tumor; border gradient is taken over the
one-pixel dilation minus erosion of the mask, and is intended to capture how
abruptly the lesion terminates.

\emph{Size.} Tumor pixel count, included explicitly so that its contribution is
visible and can be controlled for, as Section~\ref{sec:size} does.

\subsection{The Feature Set}
$C$ alone is insufficient, for reasons the results make explicit. It is combined
with fifteen further descriptors, listed in Table~\ref{tab:features}, each chosen
because it measures something the others do not.

\begin{table}[t]
\centering
\caption{The sixteen features. Each is a quantity that can be computed by hand
from the slice and the mask.}
\label{tab:features}
\begin{tabular}{lll}
\toprule
Group & Features & Measures \\
\midrule
Fractal & $C$, $F_b$ & structural irregularity \\
Texture & $\text{GLCM}_t$, $\text{GLCM}_b$, $\rho_t$ & intensity co-occurrence \\
Entropy & $H_t$ & intensity heterogeneity \\
Lacunarity & $\Lambda_t$, $\Lambda_b$ & gappiness at scale \\
Shape & solidity, eccentricity, & outline regularity \\
        & extent, circularity & \\
Edge & density, mean $\|\nabla\|$, & boundary sharpness \\
     & border $\|\nabla\|$ & \\
Size & tumor area & extent \\
\bottomrule
\end{tabular}
\end{table}

The groups are complementary in a specific sense. A fractal dimension and a
lacunarity can coincide for regions that differ visibly, because the first
measures how detail scales and the second how the region is distributed at a
fixed scale. Shape features describe the outline rather than the interior, which
is what distinguishes a well-circumscribed meningioma from an infiltrating
glioma. Edge features describe the transition at the boundary, which reflects
infiltration into surrounding tissue. Tumor area is included explicitly rather
than left implicit, so that its contribution is visible to the analyst and can be
controlled for, as Section~\ref{sec:size} does.

\subsection{Classification and Evaluation}
\label{sec:eval}
A random forest with $300$ trees is used, with maximum depth and minimum leaf
size selected by grid search. The evaluation protocol has three properties that
together prevent the inflation described in the introduction.

Folds are formed with \texttt{GroupKFold} over five splits, with the patient
identifier as the grouping variable, so that every slice from a patient falls on
the same side of every split. The inner grid search likewise uses
\texttt{GroupKFold} over the training patients, so hyperparameter selection
cannot leak across patients either. Slice predictions are aggregated to one
prediction per patient by majority vote, and the reported figure is balanced
accuracy computed over patients.

Balanced accuracy, the unweighted mean of per-class recall, is used rather than
plain accuracy because the classes are uneven: $708$ meningioma, $1426$ glioma
and $930$ pituitary slices. A classifier that predicts glioma for every input
attains $46.5\%$ plain accuracy while carrying no diagnostic information; the
same classifier attains $0.333$ balanced accuracy.

\subsection{Implementation and Cost}
Feature extraction requires one pass over each slice and is dominated by the two
box-counting estimates, which are $O(HW)$ per scale over eight scales. On the
$2740$ usable slices the full extraction takes approximately six minutes on a
single GPU-equipped Kaggle instance, with no GPU operations involved; the cost is
in the mask operations and the box counting, both of which run on the CPU. Model
fitting with the nested grid search adds approximately ten minutes.

The complete pipeline, the extracted features and the per-fold results are
available in the repository accompanying this paper, together with the seed used
for every split, so that each reported figure can be recomputed rather than taken
on trust.

%==============================================================
\section{Results}

\subsection{What the Descriptor Separates}
Fig.~\ref{fig:descriptor}(a) reports the one-vs-rest AUC of $C$ used alone as a
single-feature classifier.

\begin{figure}[t]
\centering
\includegraphics[width=\columnwidth]{fig_descriptor.png}
\caption{(a) One-vs-rest AUC of $C$ alone. The descriptor separates pituitary
adenomas well and meningiomas barely. (b) Class separation retained after
removing the component of $C$ predictable from tumor area.}
\label{fig:descriptor}
\end{figure}

The values are $0.907$ for pituitary adenoma, $0.735$ for glioma and $0.613$ for
meningioma. A single ratio therefore carries most of the information needed to
identify one class and little of what is needed to identify another.

The pattern is consistent with what the ratio measures. Pituitary adenomas are
compact and comparatively homogeneous, so the dimension of the tumor region sits
well below that of the surrounding brain and $C$ is low. Gliomas are irregular
and heterogeneous, raising $F_t$ towards or above $F_b$. Meningiomas fall between
the two and, more importantly, vary widely among themselves: some are smooth and
well circumscribed while others are lobulated with irregular margins. A single
scalar summarising irregularity cannot separate a class whose defining property
is not a consistent degree of irregularity.

We report this rather than only the aggregate model performance because it
delimits the claim. $C$ is a useful descriptor, not a sufficient one, and the
class for which it is weakest is identified.

\subsection{Is the Descriptor Measuring Size?}
\label{sec:size}
Fractal estimates are sensitive to the area of the region measured, and the three
tumor types differ in typical size, so a descriptor that appears to capture
morphology could be capturing size. The alternative explanation has to be tested
rather than dismissed.

We removed from $C$ the component predictable from tumor area and recomputed the
class separation. Fig.~\ref{fig:descriptor}(b) shows that $45\%$ of the
separation remains. The descriptor is therefore partly confounded with size and
partly independent of it. Both halves of that statement are part of the result:
a claim that $C$ measures morphology alone would be wrong, and a claim that it is
merely size in another form would also be wrong.

\subsection{Classification}
Under the protocol of Section~\ref{sec:eval}, the sixteen-feature model reaches a
patient-level balanced accuracy of $0.92$. Table~\ref{tab:compare} places this
against published results on the same dataset.

\begin{table}[t]
\centering
\caption{Published results on the Figshare dataset. The protocol column records
whether the evaluation groups by patient, since figures obtained under different
protocols are not directly comparable.}
\label{tab:compare}
\begin{tabular}{llcc}
\toprule
Method & Representation & Protocol & Acc. \\
\midrule
Cheng 2015 \cite{cheng2015} & BoW, ring partition & patient & $0.913$ \\
Ismael 2018 \cite{ismael2018} & DWT $+$ Gabor & split & $0.919$ \\
\textbf{This work} & \textbf{16 named features} & \textbf{patient} & $\mathbf{0.920}$ \\
Gumaei 2019 \cite{gumaei2019} & PCA-NGIST $+$ RELM & 5-fold & $0.942$ \\
Swati 2019 \cite{swati2019} & VGG19 fine-tuned & patient & $0.948$ \\
Deepak 2019 \cite{deepak2019} & GoogLeNet $+$ SVM & patient & $0.980$ \\
\bottomrule
\end{tabular}
\end{table}

\begin{figure}[t]
\centering
\includegraphics[width=0.88\columnwidth]{fig_literature.png}
\caption{Reported accuracy against the size and nature of the representation. The
present result is competitive with the benchmark published alongside the dataset
and below two later methods that use higher-dimensional or learned features.}
\label{fig:literature}
\end{figure}

The result exceeds the benchmark published with the dataset by less than one
point, which we do not treat as a meaningful margin. Two later methods are
higher: Gumaei \emph{et al.} by $2.2$ points using a hybrid feature
representation, and Deepak and Ameer by $6$ points using transferred deep
features under a patient-level protocol.

We state this plainly because the alternative invites a reader to make the
comparison for us. The claim of this paper is not that sixteen interpretable
features outperform a transferred convolutional network. It is that they come
within six points of one while remaining individually nameable, measurable and
open to inspection, which is a different property and useful for a different
reason.

%==============================================================
\section{Discussion}

Three observations follow from the results.

The descriptor behaves as its construction predicts. $C$ was defined to measure
irregularity relative to the host brain, and it separates most strongly the class
that is most consistently regular relative to its surroundings. It separates
weakly the class whose members vary most among themselves. A descriptor that
performed uniformly across all three would in fact be harder to interpret, since
its behaviour would not follow from what it measures.

The size analysis marks the boundary of the morphological claim. Retaining $45\%$
of separation after removing size means the descriptor carries information that
tumor area does not, and also that a substantial part of its apparent
discriminative power is area in another form. For a descriptor intended to be
interpretable this is worth knowing precisely, because an interpretation that
attributes the whole effect to structural irregularity would be wrong.

The comparison with the literature is a trade rather than a defeat. The gap to
the strongest reported method is six points of accuracy. What is obtained in
exchange is a representation in which every dimension has a name and a
computation: a clinician can ask why a slice was assigned to a class and receive
an answer in terms of measured lacunarity, boundary gradient and the fractal
ratio, each of which can be checked against the image. Whether that exchange is
worth making depends on the setting, and we do not claim it is always worth
making.

Finally, the setting in which this trade is most defensible is worth stating
explicitly, since a general claim would be too strong. Three conditions make a
named-feature representation preferable: the training cohort is small enough that
a learned representation is weakly constrained; a mask or region of interest is
available for other reasons, so computing region-based descriptors adds little
cost; and the prediction must be explained to someone who will act on it. All
three hold in the setting studied here. Where a large cohort is available and the
prediction is used for triage rather than justification, the argument runs the
other way, and the accuracy difference in Table~\ref{tab:compare} is the relevant
quantity.

%==============================================================
\section{Limitations}
\label{sec:limits}

\textbf{A mask is required at inference.} The descriptor is a ratio between two
regions, so those regions must be delineated before it can be computed. We used
expert masks in order to isolate the descriptor from segmentation error. A
deployed system would need an automatic region extraction step, and the accuracy
reported here would not transfer unchanged to that setting; the sensitivity of
$C$ to mask error is not characterised in this work and would need to be measured
directly.

\textbf{The descriptor is weak for meningioma.} An AUC of $0.613$ is close
enough to chance that $C$ should not be relied on for that class. The full
feature set compensates, but a reader interested in the descriptor itself should
take this as its principal limitation.

\textbf{Partial size confounding.} $45\%$ of separation survives the removal of
size, which means $55\%$ does not.

\textbf{Single dataset, single modality.} All results are on contrast-enhanced
T1-weighted slices from one collection. Fractal estimates depend on resolution
and reconstruction, so transfer to other acquisition protocols is not established
here.

\textbf{Two published methods are more accurate.} Stated in
Table~\ref{tab:compare} and repeated here so that it is not overlooked.

\textbf{The novelty claim is bounded.} Fractal analysis of brain tumors, and the
computation of dimensions over both tumorous and non-tumorous tissue, are
established. What we have not found in the literature is this particular ratio
used as a descriptor for three-class typing. A literature search establishes
absence of evidence rather than evidence of absence, and the claim is phrased
accordingly.

%==============================================================
\section{Conclusion}

A single ratio between the fractal dimension of a tumor and that of the brain
tissue surrounding it is comparable across patients, because both quantities come
from one acquisition and the factors that differ between acquisitions divide out.
Used alone it identifies pituitary adenomas well and meningiomas poorly, and
roughly half of its discriminative power is attributable to tumor size. Combined
with fifteen further named descriptors it classifies tumor type at $0.92$
balanced accuracy at the patient level, close to the benchmark published with the
dataset and below the strongest transferred deep representation. The value of the
approach lies in what can be said about a prediction rather than in the accuracy
of the prediction itself, and the analysis above is intended to make clear both
what the descriptor measures and where it stops being useful.

%==============================================================
\begin{thebibliography}{00}

\bibitem{cheng2015}
J. Cheng, W. Huang, S. Cao, R. Yang, W. Yang, Z. Yun, Z. Wang, and Q. Feng,
``Enhanced performance of brain tumor classification via tumor region
augmentation and partition,'' \emph{PLoS ONE}, vol.~10, no.~10, art.~e0140381,
2015, doi: 10.1371/journal.pone.0140381.

\bibitem{gumaei2019}
A. Gumaei, M. M. Hassan, M. R. Hassan, A. Alelaiwi, and G. Fortino, ``A hybrid
feature extraction method with regularized extreme learning machine for brain
tumor classification,'' \emph{IEEE Access}, vol.~7, pp.~36266--36273, 2019.

\bibitem{deepak2019}
S. Deepak and P. M. Ameer, ``Brain tumor classification using deep CNN features
via transfer learning,'' \emph{Comput. Biol. Med.}, vol.~111, p.~103345, 2019.

\bibitem{swati2019}
Z. N. K. Swati, Q. Zhao, M. Kabir, F. Ali, Z. Ali, S. Ahmed, and J. Lu, ``Brain
tumor classification for MR images using transfer learning and fine-tuning,''
\emph{Comput. Med. Imaging Graph.}, vol.~75, pp.~34--46, 2019,
doi: 10.1016/j.compmedimag.2019.05.001.

\bibitem{afshar2019}
P. Afshar, K. N. Plataniotis, and A. Mohammadi, ``Capsule networks for brain
tumor classification based on MRI images and coarse tumor boundaries,'' in
\emph{Proc. IEEE Int. Conf. Acoust. Speech Signal Process. (ICASSP)}, 2019,
pp.~1368--1372.

\bibitem{ismael2018}
M. R. Ismael and I. Abdel-Qader, ``Brain tumor classification via statistical
features and back-propagation neural network,'' in \emph{Proc. IEEE Int. Conf.
Electro/Information Technology (EIT)}, 2018, pp.~252--257.

\bibitem{yagis2021}
E. Yagis \emph{et al.}, ``Effect of data leakage in brain MRI classification
using 2D convolutional neural networks,'' \emph{Sci. Rep.}, vol.~11, art.~22544,
2021, doi: 10.1038/s41598-021-01681-w.

\bibitem{iftekhar2003}
K. M. Iftekharuddin, W. Jia, and R. Marsh, ``Fractal analysis of tumor in brain
MR images,'' \emph{Mach. Vis. Appl.}, vol.~13, no.~5--6, pp.~352--362, 2003.

\bibitem{frontiers2023}
D. Battalapalli, S. Vidyadharan, B. V. V. S. N. Prabhakar Rao, P. Yogeeswari,
C. Kesavadas, and V. Rajagopalan, ``Fractal dimension: analyzing its potential
as a neuroimaging biomarker for brain tumor diagnosis using machine learning,''
\emph{Front. Physiol.}, vol.~14, art.~1201617, 2023,
doi: 10.3389/fphys.2023.1201617.

\bibitem{martin2022}
J. Sánchez and M. Martín-Landrove, ``Morphological and fractal properties of
brain tumors,'' \emph{Front. Physiol.}, vol.~13, art.~878391, 2022,
doi: 10.3389/fphys.2022.878391.

\bibitem{haralick1973}
R. M. Haralick, K. Shanmugam, and I. Dinstein, ``Textural features for image
classification,'' \emph{IEEE Trans. Syst., Man, Cybern.}, vol.~SMC-3, no.~6,
pp.~610--621, 1973.

\bibitem{allain1991}
C. Allain and M. Cloitre, ``Characterizing the lacunarity of random and
deterministic fractal sets,'' \emph{Phys. Rev. A}, vol.~44, no.~6,
pp.~3552--3558, 1991.

\bibitem{breiman2001}
L. Breiman, ``Random forests,'' \emph{Mach. Learn.}, vol.~45, no.~1, pp.~5--32,
2001.

\bibitem{sarkar1994}
N. Sarkar and B. B. Chaudhuri, ``An efficient differential box-counting approach
to compute fractal dimension of image,'' \emph{IEEE Trans. Syst., Man, Cybern.},
vol.~24, no.~1, pp.~115--120, 1994.

\end{thebibliography}

\end{document}
